%% LyX 2.3.6.2 created this file.  For more info, see http://www.lyx.org/.
%% Do not edit unless you really know what you are doing.
\documentclass[11pt,english]{article}
\renewcommand{\familydefault}{\rmdefault}
\usepackage[T1]{fontenc}
\usepackage[latin9]{inputenc}
\usepackage{geometry}
\geometry{verbose,tmargin=0.75in,bmargin=0.75in,lmargin=0.75in,rmargin=0.75in}
\usepackage{fancyhdr}
\pagestyle{fancy}
\usepackage{babel}
\usepackage{array}
\usepackage{multirow}
\usepackage{amsmath}
\usepackage{amsthm}
\usepackage{amssymb}
\usepackage{setspace}
\usepackage[authoryear]{natbib}
\onehalfspacing
\usepackage[unicode=true,pdfusetitle,
 bookmarks=true,bookmarksnumbered=false,bookmarksopen=false,
 breaklinks=false,pdfborder={0 0 0},pdfborderstyle={},backref=false,colorlinks=false]
 {hyperref}

\makeatletter

%%%%%%%%%%%%%%%%%%%%%%%%%%%%%% LyX specific LaTeX commands.
%% Because html converters don't know tabularnewline
\providecommand{\tabularnewline}{\\}

%%%%%%%%%%%%%%%%%%%%%%%%%%%%%% Textclass specific LaTeX commands.
\theoremstyle{plain}
\newtheorem{prop}{\protect\propositionname}
\theoremstyle{plain}
\newtheorem{lem}{\protect\lemmaname}
\theoremstyle{plain}
\newtheorem*{cor*}{\protect\corollaryname}

\@ifundefined{date}{}{\date{}}
%%%%%%%%%%%%%%%%%%%%%%%%%%%%%% User specified LaTeX commands.
\usepackage{lscape}
\usepackage{hyperref}
\usepackage{fancyhdr}
\usepackage{bookmark}
\usepackage{array}
\usepackage{refstyle}
\usepackage{float}
\usepackage{dcolumn}
\usepackage{capt-of}
\usepackage{multirow}
\usepackage{titling}
\usepackage{pdfpages}

\definecolor{darkred}{rgb}{0.8,0,0}
\definecolor{darkred}{rgb}{0.7, 0.11, 0.11}

\hypersetup{colorlinks=true,citecolor=black, linkcolor=darkred}
\hypersetup{urlcolor=darkred}

\newref{fig}{refcmd={\hyperref[#1]{Figure \ref{#1}}}}
\newref{tab}{refcmd={\hyperref[#1]{Table \ref{#1}}}}

\newcolumntype{P}[1]{>{\centering\arraybackslash}p{#1}}
\newcolumntype{M}[1]{>{\centering\arraybackslash}m{#1}}

\newcolumntype{C}{>{\centering\arraybackslash}p{0.10\textwidth}}
\lhead{}
\chead{}
\rhead{}
\lfoot{}
\cfoot{\thepage}
\rfoot{}
\renewcommand{\headrulewidth}{0pt}

\newcommand\blfootnote[1]{%
  \begingroup
  \renewcommand\thefootnote{}\footnote{#1}%
  \addtocounter{footnote}{-1}%
  \endgroup
}

\setlength{\droptitle}{-50pt}
%%\setlength{\footskip}{20pt}

\makeatother

\providecommand{\corollaryname}{Corollary}
\providecommand{\lemmaname}{Lemma}
\providecommand{\propositionname}{Proposition}

\begin{document}
\pagebreak{}

\subsection{Theoretical Model\label{subsec:Theoretic-Model}}

Consider an environment with a two-tiered production network. Downstream
firms, indexed by $j$, sell differentiated varieties to consumers.
Upstream firms, indexed by $f$, sell inputs used in production by
the downstream firms. Firms, both upstream and downstream, will be
able to alter their productivity through their choice of R\&D, represented
by $I$, or production structure, represented by $\omega$. Aggregate
R\&D, $\boldsymbol{\mathcal{I}}$, and aggregate network structure,
$\boldsymbol{\Omega}$, will depend on firm-level choices and additionally
influence production technology. Factor prices are determined by a
freely-traded outside good.\footnote{I interpret the model results as characterizing optimal policy for
one sector, sufficiently small that outcome on this sector do not
affect economy-wide aggregates like factor prices or preferences.} 

\paragraph{Technology. }

I specify technology in terms of a upstream and downstream firm's
 cost of production:
\[
\begin{alignedat}{2}\text{\ensuremath{c_{f,j}\Bigl(\boldsymbol{I}_{f}},\ensuremath{\boldsymbol{\omega}_{f}},\ensuremath{\boldsymbol{\mathcal{I}}},\ensuremath{\boldsymbol{\Omega}\Bigr)}} & \text{ and } & c_{j}\Bigl(\{\bar{p}_{f,j}\}_{f},\boldsymbol{I}_{j},\boldsymbol{\omega}_{j},\boldsymbol{\mathcal{I}},\boldsymbol{\Omega}\Bigr)\end{alignedat}
.
\]
Each firm can affect productivity through R\&D expenditures $\{\boldsymbol{I}_{j},\boldsymbol{I}_{f}\}$
- potentially along a number of input-specific, buyer-specific, or
country-specific dimensions. Each firm chooses a production structure
that alters costs, $\{\boldsymbol{\omega}_{j},\boldsymbol{\omega}_{f}\}$.
As an example, this may represent locations decisions - both the choice
of where to produce and where to do R\&D - or firm-to-firm linkages
made between suppliers and buyers. Downstream costs additionally depend
on a vector of input prices, $\bar{p}_{f,j}$, provided by upstream
firms.\footnote{The dependence of costs on factor prices is suppressed into the choice
of production structure $\omega_{j},\omega_{f}$. For instance, $\omega$
may denote the location of production. If locations vary only in the
cost of labor, differences in cost under various values of $\omega$
simply reflects difference in wages.} 

The aggregates states include total R\&D expenditure, $\boldsymbol{\mathcal{I}}$,
and industry structure, $\boldsymbol{\Omega}$, additionally affect
costs. The state variable $\boldsymbol{\mathcal{I}}$ represents a
standard form of Marshallian externalities from collective knowledge
production, while aggregate industry structure capture other forms
of spillovers like geographic agglomeration forces. These aggregates
can be multi-dimensional and depend differently on the components
of R\&D and production structure.\footnote{For example, there may be country-specific aggregate R\&D levels that
depend only on R\&D spending done in those countries.}

\paragraph{Policy. }

A planner will have the opportunity to affect the prices and costs
enter the firms' problems. In particular, I write final prices $\{\bar{p}_{j}\}$,
input prices $\{\bar{p}_{f,j}\}$, and R\&D costs $\{\bar{r}_{j,k},\bar{r}_{f,l}\}$
as a function of the underlying variables and policy wedges: 
\[
\begin{alignedat}{1}\bar{p}_{j}= & \left(1-t_{j}\right)p_{j},\\
\bar{p}_{f,j}= & \left(1-t_{f}\right)\left(1-t_{f,j}\right)p_{f,j},\\
\bar{r}_{j,k}= & \left(1-s_{j,k}\right)r,\\
\bar{r}_{f,l}= & \left(1-s_{f,l}\right)r.
\end{alignedat}
\]
The indices $k$ and $\ell$ are for each separate dimension of $\boldsymbol{I}_{j}$
and $\boldsymbol{I}_{f}$.\footnote{For example, $\boldsymbol{I}_{j}$ and $\boldsymbol{I}_{f}$ may consist
of location-specific R\&D components that differentially affect cost
according to xxx.} In addition, I assume that the cost of a particular production structure
takes the following form 
\[
\bar{f}\left(\omega_{f}\right)=f\left(\omega_{f}\right)+\phi\left(\omega_{f}\right)\text{ and }\bar{f}\left(\omega_{j}\right)=f\left(\omega_{j}\right)+\phi\left(\omega_{j}\right)
\]
where $\{\phi^{U}(\omega_{j}),\phi^{D}(\omega_{j})\}$is a tax/subsidy
associated with production structure $\omega_{j}$. 

\paragraph{Profit Maximization. }

Downstream firms determine final prices, R\&D expenditures, and firm
structure to maximize profits:
\begin{equation}
\pi_{j}=\max_{\bar{p}_{j},\boldsymbol{I}_{j},\omega_{j}}q_{j}\left(p_{j}\right)\left(\bar{p}_{j}-c_{j}\left(\left\{ p_{f,j}\right\} _{f},\boldsymbol{I_{j}},\omega_{j},\boldsymbol{\mathcal{I}},\boldsymbol{\Omega}\right)\right)-\boldsymbol{\bar{r}}_{j}\boldsymbol{I}_{j}-\bar{f}\left(\omega_{j}\right).\label{eq: Downstream Profits}
\end{equation}
Downstream firms each face a residual demand curve and wield some
monopoly power. I assume that R\&D resources are in perfectly elastic
supply at a cost $\boldsymbol{\bar{r}}_{j}$ to the firm. Finally,
each production structure is associated with fixed cost payment. 

Likewise, upstream firms face a similar profit maximization problem:
\begin{equation}
\pi_{f}=\max_{\left\{ \bar{p}_{f,j}\right\} _{j},\boldsymbol{I}_{f},\omega_{f}}\sum_{j}q_{f,j}\left(p_{f,j}\right)\left[\bar{p}_{f,j}-c_{f,j}\Bigl(\boldsymbol{I}_{f},\omega_{f},\boldsymbol{\mathcal{I}},\boldsymbol{\Omega}\Bigr)\right]-\boldsymbol{\bar{r}}_{f}\boldsymbol{I}_{f}-\bar{f}\left(\omega_{f}\right)\label{eq: Upstream Profits}
\end{equation}
where the substantive difference of note is that profits come from
the sum of sales across all buyers. 

\paragraph{Aggregate State. }

The firm-level choices of R\&D expenditure $\{\boldsymbol{I}_{f},\boldsymbol{I}_{j}\}$
and structure $\{\omega_{f},\omega_{j}\}$ determine the aggregate
R\&D levels and industry structure. And these aggregate can be multi-dimensional
functions of the firms' choices:

\begin{equation}
\boldsymbol{\mathcal{I}}=\boldsymbol{\mathcal{I}}\left(\{\boldsymbol{I}_{f},\boldsymbol{I}_{j}\}\right)\text{ and }\boldsymbol{\Omega}=\boldsymbol{\Omega}\left(\{\omega_{f},\omega_{j}\}\right).\label{eq: Aggregation}
\end{equation}


\paragraph{Utility Maximization. }

A representative agent chooses an optimal consumption profile over
the varieties produces by downstream firms and an outside good $C$
\begin{equation}
U=\max_{C,\left\{ c_{j}\right\} _{j}}u\left(C,\left\{ c_{j}\right\} _{j}\right)\label{eq: Utility}
\end{equation}
subject to $C+\sum p_{j}c_{j}\leq\boldsymbol{w}\boldsymbol{L}^{\prime}+\sum_{f}\pi_{f}+\sum_{j}\pi_{j}+T.$
I set as numeraire the price of the outside good $C$, use $\boldsymbol{w}\boldsymbol{L}^{\prime}$
to denote the total payments received by factors (including payments
for R\&D and entry), and $T$ to denote the net government revenue
from taxes and subsidies. 

\paragraph{Equilibrium.}

An equilibrium is defined by a set of prices, R\&D investments, and
production structures that solve \ref{eq: Downstream Profits}, \ref{eq: Upstream Profits},
and \ref{eq: Utility} - subject to \ref{eq: Aggregation} as well
as standard market clearing and budget constraints. Note, this model
may admits the possibility of multiple equilibria depending on the
extent of complementarities in R\&D and network spillovers. Policy
has the opportunity to improve outcomes either by fixing allocations
within a given equilibrium, or by directing firms towards selecting
an equilibrium with higher aggregate welfare. 

\subsubsection{Optimal Policy}

A planner's objective is to choose a set of policy instruments, $\{t_{j},t_{f},t_{f,j},s_{j,k},s_{f,l},\phi(\omega_{f}),\phi(\omega_{j})\}$,
that maximize welfare according to \ref{eq: Utility}. I characterize
some features of the optimal industry policy
\begin{prop}
\label{prop:Optimal-Tax-1}Assume the choice of production structure
is fixed. A planner chooses a standard set of taxes on final production
and supplier production that corrects the monopoly distortion. They
choose a set of taxes on the purchase of inputs that satisfies
\[
1-t_{f,j}=\frac{q_{j}}{q_{f,j}}\frac{\partial c_{j}}{\partial p_{f,j}}
\]
and they choose an R\&D subsidy that satisfies 
\[
s_{f}=\sum_{j}\frac{\partial c_{j}}{\partial\boldsymbol{\mathcal{I}}}\frac{\partial\boldsymbol{\mathcal{I}}}{\partial\boldsymbol{I}_{f}}q_{j}+\sum_{f,j}\frac{\partial c_{f,j}}{\partial\boldsymbol{\mathcal{I}}}\frac{\partial\boldsymbol{\mathcal{I}}}{\partial\boldsymbol{I}_{f}}q_{f,j}.
\]
\end{prop}
The first tax corrects the extent to which upstream firms do not fully
internalizes the benefits of price-reductions to downstream firms.
If Shephard's lemma held for the cost function implied by \ref{eq: Downstream Profits},
then the first tax would not be necessary.\footnote{Solving \ref{eq: Downstream Profits} for the optimal allocation of
R\&D implies a cost-function that can be written as a function of
factor prices and quantities. See Appendix Section x for an example.
If Shephard's lemma held, then $q_{j}\frac{\partial c_{j}}{\partial p_{f,j}}=q_{f,j}$
and the implied tax is $0$. } However, in the presence of an endogenous R\&D allocation and thus
potential non-constant returns to scale, a reduction in the price
of inputs can lead downstream firms to change their R\&D expenditure
and affect their unit cost - an effect no internalized by suppliers.
The second tax corrects for the extent to which all firms do not fully
internalize how R\&D investment spills over into the aggregate R\&D
state $\boldsymbol{\mathcal{I}}$. 

The next set of policies regards correcting the choice of $\omega_{j}$.
I can not characterize a fully generic policy here, so I specialize
the model to a simplified case with only two types of upstream and
downstream firms. As an intuitive example, one can think of these
choices are representing the choice of suppliers or downstream firms
to offshore production.

To make sense of what policy may look like for production structures,
consider the following case. There is a unit mass of homogenous upstream
and downstream firms. Firms can choose one of two structures $\omega_{f}\in\left\{ 0,1\right\} $
and $\omega_{j}\in\left\{ 0,1\right\} $ such that the mass of type
1 upstream and downstream firms can be represented by $M$ and $N$.
Suppose that the cost function is monotonic in all arguments. One
can heuristically think about these as extensive margin choices to
offshore production. 
\begin{prop}
\label{prop:Optimal-Tax-2}Consider the case described in the previous
paragraph. Suppose an equilibrium allocation is interior, so that
$M,N\in\left(0,1\right)$, and marginal changes in M and N are well
defined. Then, the optimal tax on type $1$ firms is
\[
\left(\sum_{j}\frac{\partial c_{j}}{\partial\boldsymbol{\Omega}}q_{j}+\sum_{f,j}\frac{\partial c_{f,j}}{\partial\boldsymbol{\Omega}}q_{f,j}\right)\times\frac{d\Omega}{dM}+\left(\sum_{j}\frac{\partial c_{j}}{\partial\boldsymbol{\mathcal{I}}}q_{j}+\sum_{f,j}\frac{\partial c_{f,j}}{\partial\boldsymbol{\mathcal{I}}}q_{f,j}\right)\times\left(\boldsymbol{I}_{f|\omega_{f}=1}-\boldsymbol{I}_{f|\omega_{f}=0}\right)=\phi_{f}.
\]
 
\end{prop}
A tax is required to correct two externalities through the aggregate
states. First, firms do not internalize how their choice to offshore
affects the aggregate industry structure which alters the unit cost
of other firms. For instance, if there are geographic agglomeration
affects from the proximity of manufacturing. Second, firms do no internalize
that the different R\&D levels that come with offshoring will affect
other firms through the Marshallian externalities. For example, if
offshoring is a substitute for R\&D or causes firms to relocate R\&D
abroad.\footnote{I model R\&D as a productivity-improving endeavor. As a result, the
marginal returns to R\&D are high when costs are high. However, high
costs also tend to lower quantities sold, and the marginal returns
to R\&D rise with quantities. In order for offshoring to be a substitute
for R\&D, the cost reduction has to offset the quantity gains. }

{[}X{]} Second-Best Policy Instruments, {[}X{]} Clean Proofs, {[}X{]}
Strategic Considerations: prop4\_offshoring\_pd.lyx

\pagebreak{}

\subsection{Parametric Model\label{subsec:Parametric-Model}}

In order to take the model to the data, formalize additional theoretical
results, and make precise quantitative statements, I need to parametrize
the model. This involves taking stances on the functional forms for
costs and preferences. This section implements a simplified model
for applied work, but in the later sections I enrich the model further
for an empirical implementation. 

\paragraph{Technology. }

I start by specifying the cost of production for downstream firms:

\begin{equation}
c_{j}=\frac{\bar{c}_{j,A}}{I_{j,A}^{\eta}}+\frac{p_{j,I}}{I_{j,I}^{\eta}}\label{eq:Downstream-Costs}
\end{equation}
where $\bar{c}_{j,A}$ is the cost of assembly, $p_{j,I}$ is the
price of inputs, and $I_{j,A},I_{j,I}$ represent R\&D on the components
of production.\footnote{An additive cost function implies that assembly and inputs are perfect
complements in the downstream production function. In the empirical
implementation I allow for imperfect substitution in the downstream
production function, and estimation suggests inputs are complements.
Derivations for a CES downstream production function between assembly
and inputs are available in the Appendix.} Total R\&D is a aggregation of country-level R\&D expenditures, where
countries differ in their R\&D productivity either due to the stock
of local knowledge $\mathcal{I}_{c}$ or manufacturing proximity effects:
\begin{equation}
I_{j,A}=\left(\sum_{c\in C_{j}}\left[I_{j,A,c}\zeta_{f,A,c}\mathcal{I}_{c}^{\omega}d_{j,A,c}^{\tau_{j}}\right]^{\varsigma}\right)^{1/\varsigma}\text{ and }I_{j,I}=\ensuremath{\left(\sum_{c\in C_{j}}\left[I_{j,I,c}\zeta_{f,I,c}\mathcal{I}_{c}^{\omega}d_{j,I,c}^{\tau_{jf}}\right]^{\zeta}\right)^{1/\varsigma}}\label{eq:Downstream-R=000026D}
\end{equation}
where $d_{j,A,c}$ is the proximity of firm $j$ to its own assembly
plants while $d_{j,I,c}$ is the proximity to its suppliers manufacturing
plants. The summand is over the set of countries $C_{j}$ that firm
$j$ operates R\&D facilities in. Finally, the price of inputs is
given by 
\[
p_{j,I}^{1-\theta}=\sum_{f\in M_{j}}\left(p_{f,j}\right)^{1-\theta}
\]
where $M_{j}$ is the set of supplier faced by $j$.\footnote{There are several ways to micro-found input demand of this firm. One
interpretation is OEMs have a CES aggregation of supplier-specific
varieties. The demand system for inputs specified in \ref{eq:OEM-offer}
follows from such a setup. } 

Next, for upstream firms, the cost of production is given by
\[
c_{f,j}=\frac{\bar{c}_{f,j}}{I_{f}^{\rho}I_{f,j}^{\rho}}
\]
where $\bar{c}_{f,j}$ is the baseline cost of production by supplier
$f$ for buyer $j$ and $I_{f},I_{f,j}$ represent a generic and buyer-specific
(specialized) component of R\&D. As before, R\&D is a aggregation
of country-level R\&D expenditures where countries differ in their
R\&D productivity either due to the stock of local knowledge $\mathcal{I}_{c}$
or manufacturing proximity effects :
\begin{equation}
I_{f}=\left(\sum_{c\in C_{f}}\left[I_{f,c}\zeta_{f,c}\mathcal{I}_{c}^{\omega}d_{f,c}^{\tau_{f}}\right]^{\varsigma}\right)^{1/\varsigma}\text{ and }I_{j,I}=\ensuremath{\left(\sum_{c\in C_{f}}\left[I_{f,j,c}\zeta_{f,j,c}\mathcal{I}_{c}^{\omega}d_{f,j,c}^{\tau_{fj}}\right]^{\zeta}\right)^{1/\varsigma}}\label{eq:Upstream-R=000026D}
\end{equation}
where $d_{f,c}$ is the proximity of firm $f$ to its own manufacturing
plants while $d_{f,j,c}$ is the proximity to its buyers manufacturing
plants. The summand is over the set of countries $C_{f}$ that firm
$f$ operates R\&D facilities in.

In the notation of the theoretical model, the firm's structure can
be summarized by the tuple $\boldsymbol{\omega}_{j}\equiv\{\bar{c}_{j,A},\{d_{j,A,c}\}_{c},C_{j},M_{j}\}$
and for suppliers $\boldsymbol{\omega}_{f}\equiv\{\bar{c}_{f,j},\{d_{j,f}\}_{c},C_{f}\}$.
A firms production structure determines their baseline production
costs, their proximity to their own plants, their set of R\&D facilities,
and their set of suppliers. The R\&D vectors are summarized by $\boldsymbol{I}_{j}\equiv\{I_{j,A,c},I_{j,I,c}\}_{c}$
and $\boldsymbol{I}_{f}\equiv\{I_{f,c},\{I_{f,j,c}\}_{f}\}_{c}$ -
the country-level expenditures on assembly and input R\&D as well
as generic and specialized R\&D. And finally, the aggregate states
includes local R\&D expenditure $\boldsymbol{\mathcal{I}}\equiv\{\mathcal{I}_{c}^{\omega}\}_{c}$
and the proximity to partner firms $\boldsymbol{\Omega}\equiv\left\{ d_{j,I,c},d_{f,j,c}\right\} .$
XXX {[}partner plant locations determine part of $\bar{c}_{f,j}${]}

\paragraph{Utility Maximization. }

In the baseline model, I consider CES demand for the differentiated
varieties produced by downstream firms, subject to a fixed budget
constraint:

\[
\max_{x_{j}}\left(\sum_{j\in N}x_{j}^{\frac{\sigma-1}{\sigma}}\right)^{\frac{\sigma}{\sigma-1}}\text{ subject to }\sum_{j\in N}x_{j}p_{j}\leq\bar{D}.
\]
But in the Appendix Section X, I consider an extension to quasi-linear
preferences where the aggregate industry price index determines the
total level of demand. 

\paragraph{Profit Maximization. }

Firm optimization unfolds over two stages. First, firms determine
their production structure $\omega_{j},\omega_{f}$. Second, R\&D
is allocated, input and final prices are determined, and goods are
produced and consumed. I start by characterizing the second stage.
In particular, the downstream firms choose final prices and R\&D:
\begin{equation}
\pi_{j}\left(\boldsymbol{\omega}_{j},\boldsymbol{\Omega}\right)=\max_{p_{j},\{I_{j,A,c},I_{j,I,c}\}_{c}}D\left(p_{j}\right)^{-\sigma}\left(p_{j}-c_{j}\left(\{p_{f,j}\}_{f},\{I_{j,A,c},I_{j,I,c}\},\boldsymbol{\omega}_{j},\boldsymbol{\mathcal{I}},\boldsymbol{\Omega}\right)\right)-\sum_{c\in C_{j}}r_{c}\Bigl(I_{j,A,c}+I_{j,I,c}\Bigr)\label{eq:Upstream-Profits-Param}
\end{equation}
where $D=\bar{D}P^{\sigma-1}$ and $P^{1-\sigma}=\sum_{j}p_{j}^{1-\sigma}$.
Then, upstream firms determine R\&D. 
\begin{equation}
\pi_{j}\left(\boldsymbol{\omega}_{j},\boldsymbol{\Omega}\right)=\max_{\{I_{f,c},\{I_{f,j,c}\}_{f}\}_{c}}\sum_{j}q_{f,j}\left(p_{f,j}-c_{f,j}\left(\{I_{f,c},I_{f,j,c}\},\boldsymbol{\omega}_{f},\boldsymbol{\mathcal{I}},\boldsymbol{\Omega}\right)\right)-\sum_{c\in C_{f}}r_{c}\Bigl(I_{f}+\sum_{j}I_{f,j}\Bigr)\label{eq:Downstream-Profits-Param}
\end{equation}
And finally, input prices are determined simultaneously by OEM's take
it or leave it offer. OEM's, utilizing their larger market position,
choose input prices that minimize their input expenditure subject
to the participation of their supplier: 
\begin{equation}
p_{f,j}=\arg\min_{p_{f,j}}\underbrace{D\left(p_{j}\right)^{-\sigma}\frac{p_{j,I}}{I_{j,I}^{\alpha}}\frac{p_{f,j}^{-\theta}}{p_{j,I}^{1-\theta}}}_{\equiv q_{f,j}}p_{f,j}\text{ \text{subject to }}\tilde{\pi}_{f,j}>0.\label{eq:OEM-offer}
\end{equation}
The condition $\tilde{\pi}_{f,j}>0$ is a participation constraint
that ensures that supplier $f$ makes positive profits from supplier
$j$ accounting for $j's$ contribute to the R\&D expenses.\footnote{The explicit participation constrain is listed in Appendix Section
X. It requires that OEMs compensate suppliers for the entire specialized
R\&D expenses and their relative share of the generic R\&D expenditure.}

The stage-1 problem is then to choose the production structure that
maximizes the stage 2 profits: 
\[
\max_{\boldsymbol{\omega}_{j}}\pi_{j}\left(\boldsymbol{\omega}_{j},\boldsymbol{\Omega}\right)-f\left(\boldsymbol{\omega}_{j}\right)\text{ and \ensuremath{\max_{\boldsymbol{\omega}_{f}}}\ensuremath{\pi_{f}\left(\boldsymbol{\omega}_{f},\boldsymbol{\Omega}\right)}}-f\left(\boldsymbol{\omega}_{j}\right).
\]
I do not characterize the optimal choices at this point. The estimation
will condition on the observes production structure, and counterfactuals
will consider exogenous changes in production structure.\footnote{In the future, I plan to solve this problem and compute counterfactuals.
I additionally want to consider strategic behavior by firms at the
point where production structure are determined. That is, compute
firms' best-responses to their competitors strategies, anticipating
the effect of all production structure on the aggregate state:

\[
\max_{\omega_{j}}\pi_{j}\left(\omega_{j},\boldsymbol{\Omega}\left(\left\{ \omega_{-j},\omega_{f}\right\} \right)\right)-f\left(\omega_{j}\right)\text{ and \ensuremath{\max_{\omega_{f}}}\ensuremath{\pi_{f}\left(\omega_{f},\boldsymbol{\Omega}\left(\left\{  \omega_{j},\omega_{-f}\right\}  \right) \right)}}-f\left(\omega_{j}\right).
\]
}

\paragraph{Aggregate State. }

The aggregator function $\boldsymbol{\Omega}$ determines manufacturing
proximity from firms' production structures, $\boldsymbol{\Omega}:\left\{ \omega_{f},\omega_{j}\right\} \mapsto d_{j,I,c},d_{f,j,c}$.
And, the aggregator $\boldsymbol{\mathcal{I}}$determines local R\&D
stocks from firms' R\&D expenditures: 
\[
\mathcal{I}_{c}=\sum_{f}\left(I_{f,c}+\sum_{j}I_{f,j,c}\right)+\sum_{j}\left(I_{j,I,c}+I_{j,A,c}\right).
\]

\pagebreak{}

\subsection{Equilibrium Allocations\label{subsec:Equilibrium-Allocations}}

In this section I characterized the allocations implied by the solutions
to \ref{eq:Upstream-Profits-Param}, \ref{eq:Downstream-Profits-Param},
and \ref{eq:OEM-offer}, starting with the optimal input contract. 
\begin{lem}
\label{lem:Optimal-Contract}The optimal contract offered by a buyer
agrees to purchase inputs at a fixed markup

\[
\mu_{f,j}=1+2\rho\alpha,
\]
or, in terms of the optimal price level, a cost-plus contract
\[
p_{f,j}=\text{Unit Cost}+\alpha\times\frac{\text{R\&D Costs}}{\text{Quantities}}.
\]
\end{lem}
The optimal contract is a cost-plus contract that reimburses the suppliers
for a multiple $\alpha$ of their R\&D investments.\footnote{I interpret $\alpha$ as reduced-form parameter that captures the
bargaining power of suppliers. When $\alpha$ is higher, suppliers
can demand greater returns to their R\&D investments. } The contracted markup is increasing in $\rho$. When the returns
to R\&D are high, OEMs allow suppliers to receive greater profits
and this incentivize innovation. Additionally, final prices $p_{j}$,
will be a standard CES markup over marginal costs. 

Next, to characterize the R\&D allocations, it will be useful to define
R\&D productivity indices. These indices come from solving out the
optimal R\&D allocation across countries, allowing us to compactly
write other allocation results. For example, with supplier's generic
R\&D, we use \ref{eq:Upstream-R=000026D} to write
\begin{equation}
\bar{\zeta}_{f}=\left(\sum_{c\in C_{f}}\left[\zeta_{f,c}\mathcal{I}_{c}^{\omega}d_{f,c}^{\tau_{f}}r_{c}^{-\zeta}\right]^{\frac{\varsigma}{1-\varsigma}}\right)^{\frac{1-\varsigma}{\varsigma}}.\label{eq:Productivity-Index}
\end{equation}
This index embeds the optimal allocation of R\&D across countries,
and effectively summarizes the marginal cost of raising the level
of $I_{f}$. Appendix Section X characterizes an equivalent index
for specialized R\&D $\bar{\zeta}_{f,j}$ and OEM's R\&D $\bar{\zeta}_{j,I},\bar{\zeta}_{j,A}$
using the expressions in \ref{eq:Downstream-R=000026D} and \ref{eq:Upstream-R=000026D}.
With this notation defined, I present the following allocation results: 
\begin{prop}
\label{prop: Suppliers R=000026D Allocation-1}Upstream firms allocate
their R\&D across generic and specialized R\&D, and within specialized
R\&D, as follows: 

\[
\frac{I_{f}}{\bar{\zeta}_{f}}=\sum_{j}\frac{I_{f,j}}{\bar{\zeta}_{f,j}}\text{ and }\frac{I_{f,j}}{\sum I_{f,j}}=\frac{\left(\bar{\zeta}_{f,j}q_{f,j}\bar{c}_{f,j}\right)^{\frac{1}{1+\rho}}}{\sum_{j}\left(\bar{\zeta}_{f,j}q_{f,j}\bar{c}_{f,j}\right)^{\frac{1}{1+\rho}}}.
\]
Downstream firms allocate their R\&D across components of production
according to 
\[
\frac{I_{j,I}}{I_{j,A}+I_{j,I}}=\frac{\left(\bar{\zeta}_{j,I}p_{j,I}\right)^{\frac{1}{1+\eta}}}{\left(\bar{\zeta}_{j,A}\bar{c}_{j,A}\right)^{\frac{1}{1+\eta}}+\left(\bar{\zeta}_{j,I}p_{j,I}\right)^{\frac{1}{1+\eta}}}
\]
An upstream firms allocates generic R\&D expenditures across countries
according to 
\[
\frac{I_{f,c}}{\sum_{c'}I_{f,c'}}=\frac{\left(I_{f,c}\zeta_{f,c}\mathcal{I}_{c}^{\omega}d_{f,c}^{\tau_{f}}r_{c}^{-1}\right)^{\frac{\varsigma}{1-\varsigma}}}{\sum_{c'}\left(I_{f,c'}\zeta_{f,c'}\mathcal{I}_{c'}^{\omega}d_{f,c'}^{\tau_{f}}r_{c'}^{-1}\right)^{\frac{\varsigma}{1-\varsigma}}}
\]
with an equivalent expression for specialized R\&D $I_{f,j,c}$ and
OEM's allocations $I_{j,I,c},I_{j,A,c}$. 
\end{prop}
This characterization highlights four economic forces. First, the
division between generic and specialized R\&D reflects their relative
productivity and returns. When generic R\&D is highly productive (high
$\bar{\zeta}_{f}$), suppliers tilt investment towards cost reductions
that benefits all buyers. When buyer-specific productivity is high,
specialized R\&D becomes more attractive.

Second, specialized R\&D is allocated towards buyers for which suppliers
are productive $\bar{\zeta}_{f,j}$, have high quantity $q_{f,j}$,
and have high cost $\bar{c}_{f,j}$. Higher quantities increase the
scale over which cost-reducing innovations apply, and higher cost
increase the marginal benefit of innovation.\footnote{In equilibrium, higher costs will also reduce the quantities that
suppliers receive. Both through competition in input markets, and
reduced input purchase from downstream firms. The equilibrium implications
of cost shocks should therefore not be interpreted as shocks that
necessarily increase innovation.} When the elasticity $\rho$ increases, the returns to R\&D become
stronger ad more uniform across buyers.

Third, OEMs allocate R\&D towards the components of production where
productivity improvements have the highest return - either due to
high R\&D productivity ($\bar{\zeta}_{j,I},\bar{\zeta}_{j,A}$) or
high production costs ($p_{j,I},$$\bar{c}_{j,A}$).\footnote{This result hinges on the assumption that inputs are complements.
If inputs were sufficiently substitutable, it would be the case that
more R\&D would be allocated to the components of production that
had the lowest cost. The data supports the fact that inputs are complements.}

Fourth, the cross-country allocation of R\&D reflects a standard CES
tradeoff between productivity and cost. R\&D flows to locations where
productivity is high relative to costs. And as $\varsigma\rightarrow1$,
expenditure concentrates in the single most productive location.\footnote{This theory suggests why directly interpreting reduced-form coefficients
from innovation outcomes as causal coefficients, even when using instruments,
is problematic. As R\&D will endogenously be allocated to the locations
where it is most productive.} 

\pagebreak{}

\subsection{Equilibrium Comparative Statics - First Order\label{subsec:Comparative-Statics-First-Order}}

To get a better sense of the model mechanism and the potential implications
of policy, I characterize the response of industry outcomes GDP to
price and R\&D efficiency shocks. But changes in prices and R\&D efficiency
depend on more primitive parameters, and hence we provide some results
to think about how prices respond to cost shocks (which affect productivity
by raising unit costs) and R\&D productivity shocks (which affect
the efficiency of R\&D investment). Throughout, these comparative
statics allow R\&D allocations and prices to adjust, but treat $\omega_{f},\omega_{j}$
as exogenous variables from Stage 1.

\begin{table}
\centering{}\caption{Variable Summary\label{tab:Lambda-Summary}}
\def\arraystretch{1.4}%%
\begin{tabular}{cll}
\hline 
\multicolumn{1}{c}{Variables} &  & \multicolumn{1}{c}{Description}\tabularnewline
\hline 
\hline 
\multirow{1}{*}{$\Lambda_{j}$} & $:$ & Share of $j$ in Final Expenditure\tabularnewline
$\lambda_{j,A},\lambda_{j,I}$ & $:$ & Share of Assembly/Inputs in $j$'s Cost \tabularnewline
$\lambda_{f,j}$ & $:$ & Share of Buyer $j$ in Supplier $f$'s Sales\tabularnewline
$\lambda_{j,f}$ & $:$ & Share of Supplier $f$ in Buyer $j$'s Cost\tabularnewline
$\lambda_{f,c},\lambda_{j,c}$ & $:$ & Share of Country $c$ in Firm $f$'s/$j$'s R\&D\tabularnewline
$\lambda_{c,f},\lambda_{c,j}$ & $:$ & Share of Firm $f$'s/$j$'s in Country $c$'s Total R\&D\tabularnewline
\hline 
\end{tabular}
\end{table}

I start by noting that changes in real output in this industry can
be characterized by the following comparative static:
\begin{prop}
\label{prop: Network Response-1-1}The first-order change in real
industry output is given by
\[
d\log\text{GDP}=-\sum_{j}\Lambda_{j}d\log p_{j}=-\sum_{j}\Lambda_{j}\left(\lambda_{j,A}\left(d\log\bar{c}_{j,A}-\eta\bar{\zeta}_{j,A}\right)-\lambda_{j,I}\left(d\log p_{j,I}-\eta\bar{\zeta}_{j,I}\right)\right),
\]
with $\Lambda_{j}$ measure the share of firm $j$ in final expenditure,
and $\lambda_{j,A},\lambda_{j,I}$ measure the share of assembly and
inputs in the cost of firm $j$. 
\end{prop}
Changes in real output arise from either changes in productivity,
manifesting as changes in assembly costs $\bar{c}_{j,A}$ or input
prices $p_{j,I}$, or changes in R\&D efficiency, represented by changes
in the indices $\bar{\zeta}_{j,I}$ and $\bar{\zeta}_{j,A}$. While
assembly costs are a technological parameter, inputs prices and productivity
indices are endogenous through either individual R\&D choices or aggregate
R\&D spillovers.\footnote{Assembly costs are a technological parameter as we have conditioned
on the Stage 1 choices. In principle, these too are endogenous to
the plant location and market entry choices of firms.} The following results characterize the response of each to primitive
cost or R\&D productivity shocks. 
\begin{prop}
\label{prop: Network Response}The response of input prices to a set
of buyer-supplier cost shocks $d\log\bar{c}_{f,j}$ is given by 
\[
d\log p_{f,j}=d\log\bar{c}_{f,j}-\sum_{k}\left(\delta_{2}\mathbb{1}_{j=k}+\delta_{3}\lambda_{f,k}\right)\left((1-\theta)d\log\bar{c}_{f,k}+\tilde{\lambda}_{j,I}d\log p_{k,I}\right)-\delta_{1}\sum_{k}\Lambda_{k}\lambda_{k,I}d\log p_{k,I}
\]
where $d\log p_{k,I}=\sum\lambda_{k,f}d\log p_{f,k}$. The variable
$\lambda_{f,j}$ is the share of downstream firm j in upstream firm
f's revenue, and $\lambda_{j,f}$ is the share of upstream firm $f$
in the input price of downstream firm $j$. The variable $\mathbb{1}_{j=k}$
is an indicator that takes a value of 1 when buyer $k$ in the summand
equals buyer $j$. The parameters $\delta_{1}$, $\delta_{2}$, and
$\delta_{3}$are defined below. 
\end{prop}
This comparative static reflects several forces. First, there is the
direct effect of costs to prices, which ignoring any adjustments in
quantity or R\&D, directly pass through due to the pricing implications
of Lemma \ref{lem:Optimal-Contract}. 

Second, for every buyer in a supplier's network, there is a potential
supplier-specific cost shock $d\log\bar{c}_{f,k}$, which affects
R\&D incentives as it increases costs and lowers demand (hence the
$1-\theta$). And, each buyer is also experience a change in its total
input costs $d\log p_{k,I}$, arising from changes in the prices offered
by each supplier. Total input costs affect the R\&D incentives of
suppliers through four separate mechanisms: 
\[
\tilde{\lambda}_{j,I}\equiv\underbrace{\theta-1\frac{}{\vphantom{\left(\right)}}}_{\text{Upstream Competition}}+\underbrace{1\vphantom{\frac{}{\vphantom{\left(\right)}}}}_{\text{Input Expenditure}}-\underbrace{\frac{\rho}{\rho+1\vphantom{\left(\right)}}}_{\text{Input R\&D}}-\underbrace{\frac{1}{\rho+1}\frac{\sigma}{1+\rho\left(1-\sigma\right)}\lambda_{j,I}}_{\text{Downstream Demand}}.
\]
Higher input prices faced by buyer $j$, benefit each supplier as
it perceives weaker perceive competition and higher input expenditure.
However, higher input prices lead downstream firms to reallocate R\&D
towards inputs and raise final prices, which both lower input demand.
The resulting changes in R\&D affect costs through though the reallocation
of specialized R\&D, which affects particular buyers, and expansion
of generic R\&D spending as well as the overall R\&D budget, which
benefits all buyers: 

\[
\delta_{2}\mathbb{1}_{j=k}+\delta_{3}\lambda_{f,k}=\underbrace{\frac{\rho}{1+\rho\left(1-\theta\right)}}_{\text{Reallocation of R\&D to \ensuremath{j}}}\times\mathbb{1}_{j=k}+\biggl(\underbrace{\frac{2\rho}{1-2\rho\left(\theta-1\right)}}_{\text{Generic R\&D and R\&D Budget}}-\underbrace{\frac{\rho}{1-\rho\left(\theta-1\right)}}_{\text{Reallocation of R\&D to \ensuremath{k}}}\biggr)\times\lambda_{f,k}.
\]

The final term normalizes changes in $d\log p_{k,I}$ to emphasize
that relative prices changes are what determines reallocation in downstream
demand: $\delta_{1}=\left(\delta_{2}+\delta_{3}\right)\left(\sigma-1\right)/1+\rho\left(1-\sigma\right)$.
But notably, this last term will not cancel out the other $d\log p_{k,I}$
terms since even if every firm experienced the same cost shock such
that relative quantities did not change, the absolute level of demand
would change which would affect the R\&D incentives of firms. 

Specializing Proposition \ref{prop: Network Response} to a single
buyer, removes terms associated with reallocation across buyers, and
highlights the key sources of inefficiency in the model. 
\begin{cor*}
Consider the case with only one downstream buyer, the change in real
GDP can be expressed as 

\[
d\log\text{GDP}=\frac{\lambda_{j,I}\sum_{f}d\log\bar{c}_{f,j}}{\left(1-\left(\frac{2\rho}{1+2\rho\left(1-\theta\right)}\right)\left(\frac{1}{1+\eta}\lambda_{j,I}+\frac{\eta}{1+\eta}\right)\right)}.
\]
\end{cor*}
This result echo's Hulten's theorem, but has a non-unitary coefficient.
This arises due to the presence of a baseline externality, alluded
to in Proposition \ref{prop:Optimal-Tax-1} from the failure of Shephard's
Lemma. Suppliers do not capture the full extent of their productivity
improvement. In particular, they do not capture the fact that productivity
improvements affect downstream quantities, $\frac{1}{1+\eta}\lambda_{j,I}$,
and downstream R\&D allocation, $\frac{\eta}{1+\eta}$. Next, I discuss
changes in R\&D efficiency: 
\begin{prop}
\label{prop: Network Response-1-2}Consider the response of the productivity
index $\bar{\zeta}_{f}$ to a set of shocks $\left\{ \zeta_{f,c}\right\} \forall f,c$:
\[
\begin{alignedat}{1}d\log\bar{\zeta}_{f}= & \sum_{c\in C_{f}}\lambda_{f,c}\Biggl(\underbrace{d\log\zeta_{f,c}}_{\text{Direct Productivity Effect}}+\frac{\omega}{1-\zeta}\sum_{f'}\lambda_{c,f'}\biggl(\underbrace{d\log\zeta_{f',c}-\sum_{c'\in C_{f'}}\lambda_{f,c'}d\log\zeta_{f',c'}}_{\text{Spillovers from Across-Country R\&D Reallocation}}\biggr)\\
 & \quad-\omega\biggl(\underbrace{\frac{1}{\rho}\sum_{f'}\lambda_{c,f'}\sum_{j}\lambda_{f',j}d\log p_{f',j}+\frac{1}{\eta}\sum_{j}\lambda_{c,j}d\log p_{j}}_{\text{Spillovers from R\&D Expansion}}\biggr)\Biggr)
\end{alignedat}
\]
where $\lambda_{f,c}$ is the share of R\&D expenditure by firm $f$
in country $c$, $\lambda_{c,f'}$ is the share of R\&D in country
$c$ done by firm $f'$. 
\end{prop}
This comparative static reflects several forces. First, the R\&D efficiency
index responds to R\&D productivity shocks in all countries, with
shares depending on the extent of baseline R\&D expenditure done in
each country. Second, R\&D efficiency responds to R\&D shocks at all
firms and all countries, to the extent that the strength of R\&D spillovers
measured by $\omega$ are higher. In particular, based on R\&D productivity
shocks, firms relocate their R\&D expenditure across countries - where
the higher substantially there is $\zeta$, the stronger these effects
are. If other firms are relocating R\&D to the countries where firm
$f$ does R\&D intensely, then this firm benefits. Finally, in general,
R\&D productivity changes will induce more R\&D investment by suppliers
and buyers. In the absence of cost-shocks, prices only adjust due
to R\&D changes and hence scaling prices changes by the inverse of
the returns to R\&D get the expansion in total R\&D. Obviously, prices
are a potentially complicated endogenous object, so we characterize
their response to R\&D shocks in Appendix Section X. 

\subsection{Equilibrium Comparative Statics - Second-Order}

The final result we provide in this section is a characterization
of the second-order effect of cost shocks. To parse these results,
I first introduce some notation. First, we use Proposition \ref{prop: Network Response}
to characterize the first-order price response to cost changes, which
we compactly represent as
\[
\psi_{f,k}^{j}=d\log P_{j}/d\log c_{f,k}.
\]
Second, we introduce two notions of covariance in the first-order
price response. The first of these,
\[
Cov_{\Lambda}\left(\lambda_{j,I}\psi_{f,j}^{j},\lambda_{j,I}\psi_{f,k}^{j}\right)=\sum_{j}\Lambda_{j}\lambda_{j,I}\psi_{f,k}^{j}\lambda_{j,I}\psi_{f,j}^{j}-\sum_{j}\Lambda_{j}\lambda_{j,I}\psi_{f,j}^{j}\sum_{k}\lambda_{k,I}\Lambda_{k}\psi_{f,k}^{j},
\]
measures the covariance in prices changes experienced by downstream
buyers, where we use the distribution of buyers in the downstream
price index $\Lambda_{j}$. The second of these,
\[
Cov_{\lambda_{f}}\left(\tilde{\lambda}_{j,I}\psi_{f,j}^{j},\tilde{\lambda}_{j,I}\psi_{f,k}^{j}\right)=\sum_{j}\lambda_{f,j}\tilde{\lambda}_{j,I}\psi_{f,j}^{j}\tilde{\lambda}_{j,I}\psi_{f,k}^{j}-\sum_{j}\lambda_{f,j}\tilde{\lambda}_{j,I}\psi_{f,j}^{j}\sum_{k}\lambda_{f,k}\tilde{\lambda}_{k,I}\psi_{f,k}^{l}.
\]
again measures the covariance in prices changes experienced by downstream
buyers, but instead uses the distribution of buyers in the profits
of supplier $f$. Using these definitions. We can now characterize
the response of second-order effect of cost shocks on input prices. 
\begin{prop}
\label{prop: Network Response - Second Order-1}The second-order response
of input prices to a set of buyer-supplier cost shocks \{$d\log\bar{c}_{f,j}$,$d\log\bar{c}_{f,k}$\}
where $j\neq k$ is given by 
\[
\begin{alignedat}{1} & \frac{d^{2}\log p_{f,j}}{d\log\bar{c}_{f,j}d\log\bar{c}_{f,k}}=\delta_{1}\delta_{6}Cov_{\Lambda}\left(\lambda_{j,I}\psi_{f,j}^{j},\lambda_{j,I}\psi_{f,k}^{j}\right)+\delta_{1}\delta_{7}\sum_{j}\left(\left(1-\lambda_{j,I}\right)\psi_{f,k}^{j}\right)\lambda_{j,I}\Lambda_{j}\psi_{f,j}^{j}\\
 & \quad+\delta_{2}\delta_{8}\left(1-\lambda_{j,I}\right)\psi_{f,k}^{j}\psi_{f,j}^{j}+\delta_{3}\delta_{8}\sum_{j}\lambda_{f,j}\left(1-\lambda_{j,I}\right)\psi_{f,k}^{j}\psi_{f,j}^{j}-\delta_{3}\delta_{5}Cov_{\lambda_{f}}\left(\tilde{\lambda}_{j,I}\psi_{f,j}^{j},\tilde{\lambda}_{j,I}\psi_{f,k}^{j}\right)
\end{alignedat}
.
\]
\end{prop}
Now to unpack this expression, we define each new parameter in turn.
First, recall from above that $\delta_{1}$ determines how changes
in the aggregate downstream price index affect the profit incentive
for innovating for each buyer. With that in mind, note that 
\[
\delta_{6}=\frac{\sigma-1}{1+\rho\left(1-\sigma\right)},
\]
which is the elasticity of downstream demand shares with respect to
cost changes. What this term captures is how the first-order changes
in the final prices of downstream firms relocate quantity across buyers.
These quantities in turn affect the demand faced by various suppliers
and thus their incentive to do R\&D. The covariance appears as a measure
of the extent to which buyers which are experiencing large first-order
price effects are in the second-order also experiencing reallocations
of demand shares towards them. Relatedly, we have $\delta_{7}=1/\left(1+\rho\right),$
which captures how for each downstream firm in the final price index,
the cost share of inputs is changing - and hence the ability of upstream
R\&D to influence downstream prices. 

Next, recall that $\delta_{2}$ and $\delta_{3}$ determine how changes
in the input prices of various buyers affects the R\&D incentives
of suppliers either through a competition effect or a reallocation/expansion
of input demand. Now note,
\[
\delta_{4,j}=\frac{1}{\rho+1},
\]
which similar to $\delta_{7}$, captures how within each buyer $j$
the share of costs that is determined by upstream R\&D choice is changing,
and this amplifies any first-order feedback effect between upstream
R\&D and demand for inputs. The first component of $\delta_{4,j}$
captures how this second-order effect does not alter the competition
effect from the change in input prices which is governed by $-\theta$;
it only influences the input demand expansion channel $\frac{1}{\rho+1}\frac{\sigma}{1+\rho\left(1-\sigma\right)}\lambda_{j,I}$. 

Finally, we have another change stemming from the common component
of upstream R\&D, where we have 
\[
\delta_{5}=\frac{1}{1-\rho\left(\theta-1\right)},
\]
which is the elasticity of profit shares $\lambda_{f,j}$ with respect
to profit shifters $y_{j}.$ And why this appears with a covariance
term is to measure the extent to which buyers which are experiencing
large first-order price effects are in the second-order also experiencing
reallocations of profits shares in supplier $f$ toward them.

XXX Note how this $\bar{c}_{f,j}$ can also be a wedge. Policy instrument
- effect on outcomes. 
\[
\delta_{8}=\frac{-\sigma}{1+\eta\left(1-\sigma\right)}\frac{1}{\left(1+\eta\right)^{2}}
\]

\pagebreak{}

\subsection{Specialized Equilibrium Characterization}

{[}X{]} Comparativ Static - Second Order Productivity Change

{[}X{]} Comparative Statics - Can go here next if prev is finished. 

{[}X{]} GDP Corollary
\end{document}
